\subsection{The flat-closure lemma of cote 127}

The following statement settles the existence and finite-presentation assertions
of Lemma 1.3 in cote 127, scans 3--4.  We have inspected those two scans.
The parenthetical insertion on scan 3 is the injectivity of
$A\longrightarrow\Gamma(Y',\mathcal O_{Y'})$; the prime on $Y'$ is essential.
We use ``universally separating'' in the sense specified on scan 3: after every
base change on the base, restriction to $U$ is injective as a map of sheaves.

\begin{lemma}[Known: Artinian flatness descent]\label{A4:artinianflat}
Let $A\to B$ be an injective homomorphism of commutative rings, with $A$
Artinian.  If an $A$-module $M$ has $M\otimes_A B$ flat over $B$, then $M$
is flat over $A$.  Moreover, the functor $-\otimes_A B$ detects zero modules.
\end{lemma}
\begin{proof}
The flatness assertion is Stacks, Lemma 10.101.7
\cite[Tag 051E, Lemma 10.101.7]{StacksArtinian}, whose two proofs we opened.
For the additional assertion, decompose $A$ into its finitely many Artinian
local factors; this induces the corresponding product decomposition of $B$.
Thus assume $(A,\mathfrak m,k)$ local.  The ideal $\mathfrak mB$ is nilpotent,
and $B/\mathfrak mB\ne0$: otherwise $B=0$, contradicting injectivity.
If $M\otimes_A B=0$, then
$(M/\mathfrak mM)\otimes_k(B/\mathfrak mB)=0$.  Faithfulness of tensoring
vector spaces with a nonzero vector space gives $M/\mathfrak mM=0$.
Since $\mathfrak m$ is nilpotent, this implies $M=0$ without any finite
generation hypothesis.
\end{proof}

\begin{theorem}[Proved: cote 127, Lemma 1.3, with weaker hypotheses]
\label{A4:flatclosure}
Let $A$ be Artinian, $Y=\operatorname{Spec}A$, $f:X\to Y$ any morphism of
schemes, and $j:U\hookrightarrow X$ a retrocompact open immersion.  Let
$\mathcal F$ be a quasi-coherent $\mathcal O_X$-module and let
$\mathcal F|_U\twoheadrightarrow\mathcal G_U$ be a quasi-coherent quotient
which is flat over $Y$.
Suppose there is a morphism $Y'\to Y$ with
$A\to\Gamma(Y',\mathcal O_{Y'})$ injective, and on $X'=X\times_Y Y'$ a
quasi-coherent quotient $\mathcal F_{Y'}\twoheadrightarrow\mathcal G'$ such
that $\mathcal G'|_{U'}=\mathcal G_{U'}$, $\mathcal G'$ is $Y'$-flat, and
$\mathcal G'\to j'_*\mathcal G_{U'}$ is universally injective under base
change on $Y'$.
Then
\[
 \mathcal G=\operatorname{im}(\mathcal F\longrightarrow j_*\mathcal G_U)
\]
is $Y$-flat, its restriction map is universally injective under base change
on $Y$, and its formation as this image commutes with every base change on
$Y$.  In particular $\mathcal G_{Y'}\simeq\mathcal G'$ as quotients of
$\mathcal F_{Y'}$.  It is the unique quotient having the required properties.
If $\mathcal G'$ is of finite presentation, so is $\mathcal G$.
\end{theorem}
\begin{proof}
Since $j$ is quasi-compact and quasi-separated, $j_*\mathcal G_U$ and its
image $\mathcal G$ are quasi-coherent.  All conclusions may be checked on
affine opens of $X$.  We first reduce the test morphism to an injective map
$A\to B$ with $Y'=\operatorname{Spec}B$.
Take affine opens $V_i$ covering the original $Y'$.  The intersection of
the kernels of $A\to\Gamma(V_i,\mathcal O)$ is zero.  The descending chain
condition on ideals of $A$ implies that a finite subfamily already has
zero intersection.  Its disjoint union is affine, with ring the finite
product $B$ of these rings.  The hypothesis on $\mathcal G'$ persists on
this disjoint union.  Once the asserted universal property has been proved
using this test, uniqueness gives the comparison over the original $Y'$.

Fix an affine open $W=\operatorname{Spec}R$ of $X$.  Write
$U\cap W=\bigcup_{i=1}^r D(g_i)$ using retrocompactness, and put
\[
 C=\bigoplus_{i=1}^r\Gamma(D(g_i),\mathcal G_U),\qquad
 G=\operatorname{im}(\Gamma(W,\mathcal F)\longrightarrow C),\qquad Q=C/G.
\]
Each summand of $C$ is $A$-flat because $D(g_i)$ is affine and
$\mathcal G_U$ is flat over $A$; thus $C$ is $A$-flat.  The module $G$ is
$\Gamma(W,\mathcal G)$, since restriction of a section to an open is zero
if and only if its restrictions to the chosen open cover are zero.
For any $A$-algebra $T$, the module $C\otimes_A T$ identifies canonically
with the direct sum of the sections of $\mathcal G_{U_T}$ on the affine
opens $D(g_i)_T$.

Over $B$, the image of $\Gamma(W,\mathcal F)\otimes_A B$ in $C\otimes_A B$
is precisely $G'=\Gamma(W_B,\mathcal G')$.  The map
$G'\hookrightarrow C\otimes_A B$ is universally injective as a map of
$B$-modules.  Indeed, the given universally separating property gives
injectivity after tensoring with every $B$-algebra $T$; using the square-zero
$B$-algebra $B\oplus N$ for each $B$-module $N$ gives injectivity after
tensoring with $N$ as well.  Consequently its cokernel is $B$-flat:
the Tor exact sequence, together with flatness of $C\otimes_A B$, gives
$\operatorname{Tor}_1^B(N,(C\otimes_A B)/G')=0$ for every $N$.
Right exactness identifies this cokernel with $Q\otimes_A B$.
Lemma~\ref{A4:artinianflat} therefore implies that $Q$ is $A$-flat.

The exact sequence $0\to G\to C\to Q\to0$ now remains exact after tensoring
with every $A$-module; also $G$ is $A$-flat.  Its base change to any
$A$-algebra $T$ is the image of
$\Gamma(W,\mathcal F)\otimes_A T\to C\otimes_A T$, because
$\Gamma(W,\mathcal F)\twoheadrightarrow G$ stays surjective.
This proves both universal injectivity of restriction and commutation of
the closure with base change.  These local assertions glue, and the image
description gives uniqueness.

For the finite-presentation assertion it is enough to work over $R$.
As $G\otimes_A B$ is a finite $R\otimes_A B$-module, its finitely many
generators involve finitely many elements of $G$.  Let $G_0$ be the
$R$-submodule they generate.  Then $(G/G_0)\otimes_A B=0$, whence $G=G_0$
by Lemma~\ref{A4:artinianflat}.  Choose a finite presentation start
$0\to K\to R^n\to G\to0$.  Since $G$ is $A$-flat, tensoring with $B$
is still exact on the left.  The finite-presentation hypothesis implies
that $K\otimes_A B$ is finite over $R\otimes_A B$.  The same finite-generator
and zero-detection argument shows that $K$ is finite over $R$.  Thus $G$
is finitely presented.  No finite-presentation hypothesis on $X\to Y$ was
needed.
\end{proof}

\begin{remark}[Proved: the role of the manuscript hypotheses]
\label{A4:flatclosure-scope}
The theorem applies to each of the two alternatives on cote 127, scan 4.
The extra finiteness alternatives are unnecessary for the existence claim
as long as ``universally separating'' retains its stated meaning and $U$
is retrocompact.  Merely requiring $\mathcal G'\to j'_*\mathcal G_{U'}$
to be injective would not supply the flatness of $Q\otimes_A B$ used above.
This proof does not assert effectiveness of arbitrary non-flat descent
data; it uses the specified quotient of $\mathcal F$ and its restriction to
$U$ throughout.
\end{remark}

\subsection{The square-zero boundary comparison of cote 128}

\begin{proposition}[Proved: exact correction to the connecting maps]
\label{A4:boundary}
Let $A\to A'$ be a map of commutative rings, let $I\subset A$ satisfy
$I^2=0$, and write $C^q=(A')^{\otimes_A(q+1)}$, $J^q=IC^q$ and
$\overline C^q=C^q/J^q$.  Let
$C_0=H^0(\overline C^\bullet)$ and $H=H^1(J^\bullet)$, where the
cochain complexes use the usual alternating sums of cofaces.  Then the
two kernels in cote 128, scan 20, are canonically
\[
 K_a=H/\delta(C_0),\qquad
 K_m=H/\{u^{-1}\delta(u):u\in C_0^\times\},
\]
respectively.  Here $\delta:C_0\to H$ is the additive connecting map and
the second expression is the image of a homomorphism from
$C_0^\times$ to the additive group of $H$.  In particular, the map
$1+x\mapsto x$ induces the asserted canonical identification of the
kernels if and only if the two displayed subgroups of $H$ agree.
\end{proposition}
\begin{proof}
In each degree, the square-zero ideal gives exact sequences
$0\to J^q\to C^q\to\overline C^q\to0$ and
$1\to1+J^q\to(C^q)^\times\to(\overline C^q)^\times\to1$.
The last map is surjective because units lift through a nilpotent ideal.
The degreewise maps $1+x\mapsto x$ identify $1+J^\bullet$ with
$J^\bullet$ as cosimplicial abelian groups.  The cohomology exact sequences
give the two quotient descriptions.  If $u\in C_0^\times$ and $\widetilde u$
is a unit lift in $C^0$, the multiplicative boundary is
$d^1(\widetilde u)^{-1}d^0(\widetilde u)$, which under $1+x\mapsto x$ is
$u^{-1}(d^0\widetilde u-d^1\widetilde u)$ because $J^2=0$.
This is $u^{-1}\delta(u)$.  The last assertion follows because the identity
map of $H$ descends to an isomorphism of its two quotients exactly when
their kernel subgroups agree.
\end{proof}

\begin{proposition}[Proved counterexample: the canonical comparison in characteristic two]
\label{A4:boundary-counterexample}
The hypothesis of cote 128, Lemma 1, does not make the map $1+z\mapsto z$
descend to the claimed canonical isomorphism, even for an Amitsur complex
of commutative noetherian rings.  An example is
\[
 k=\mathbf F_2,\quad R=k[x,y],\quad
 A=R[a,b]/(a^2,b^2),\quad
 I=(\epsilon),\quad\epsilon=ya+xb,\quad A'=A_x\times A_y.
\]
Here $I^2=0$, and the notation of Proposition~\ref{A4:boundary} gives
\[
 B=H^0(C^\bullet)=A,\qquad
 C_0=R[t]/(t^2),\qquad
 \operatorname{im}(B\to C_0)=R+(x,y)Rt.
\]
In particular, $C_0=C_0^\times+\operatorname{im}B$ holds.  Nevertheless,
there are independent elements $e,f\in H=H^1(J^\bullet)$ with
\[
       \delta(C_0)=k e,\qquad
       \delta_{\log}(C_0^\times)=k(e+f).
\]
Thus the two quotients in Proposition~\ref{A4:boundary} cannot be identified
by the stated square-zero map.
\end{proposition}
\begin{proof}
Put $U=D(x)\cup D(y)\subset\operatorname{Spec}A$ and
$U_R=\operatorname{Spec}R\setminus\{(x,y)\}$.
The Amitsur complex for $A_x\times A_y$ is the ordered \v Cech complex of
this two-open cover of $U$.  Its degree-zero equalizer and degree-one
cohomology may therefore be computed using the usual normalized complex
of sections on $D(x),D(y),D(xy)$.  In particular, for an $R$-module $M$
whose associated sheaf is used on $U_R$, its degree-one term is
$M_{xy}/(M_x+M_y)$.

The ring $A$ is free over $R$ with basis $1,a,b,ab$.  The identity
$R_x\cap R_y=R$ inside $R_{xy}$ gives $\Gamma(U,\mathcal O_U)=A$.
Modulo $I$, on $D(x)$ put $t=a/x$, and on $D(y)$ put $t=b/y$.
The relation $ya+xb=0$ makes these sections agree on the overlap, and
on both opens $a=xt$, $b=yt$, $t^2=0$.  Hence
$U_0\simeq U_R\times\operatorname{Spec}k[t]/(t^2)$ and
$\Gamma(U_0,\mathcal O_{U_0})=R[t]/(t^2)$.
The map from $A$ has image $R+(x,y)Rt$, since $ab$ maps to zero.
For any $r,s\in R$,
\[
             r+st=(r+1)+(1+st),
\]
where $r+1$ comes from $B$ and $1+st$ is a unit; this verifies precisely
the lemma's hypothesis, without any omitted assumption.

As an $R$-module the ideal $I$ is
\[
           R\epsilon\ \oplus\ (x,y)R\,ab.
\]
Indeed $a\epsilon=xab$, $b\epsilon=yab$, and
$ab\epsilon=0$, while the two displayed summands are independent in
the basis $a,b,ab$.  On $U_R$ the ideal sheaf is therefore free with
basis $\epsilon,ab$.  Set
\[
 W=R_{xy}/(R_x+R_y),\qquad
 H=W\epsilon\oplus W ab,\qquad
 e=\left[\frac{\epsilon}{xy}\right],\quad
 f=\left[\frac{ab}{xy}\right].
\]
The Laurent monomials $x^{-i}y^{-j}$, $i,j\geq1$, form a $k$-basis of
$W$, so $e$ and $f$ are independent and nonzero.

The lifts $a/x$ and $b/y$ of $t$ give $\delta(t)=e$.
Consequently $\delta(st)=s e=s(0,0)e$ for $s\in R$, because $x e=y e=0$.
Every element of $R$ lifts globally, so the full additive boundary image
is $k e$.  Since $R^\times=k^\times=\{1\}$, all the units of $C_0$ are
exactly $1+st$, $s\in R$.  Using the local unit lifts
$1+sa/x$ and $1+sb/y$ gives
\[
 (1+sa/x)^{-1}(1+sb/y)
   =1+\frac{s\epsilon+s^2ab}{xy},
\]
because $a^2=b^2=0$ and the characteristic is two.  Thus
\[
 \delta_{\log}(1+st)=s e+s^2 f=s(0,0)(e+f).
\]
This proves the asserted distinct boundary subgroups.

There is also a direct well-definedness witness.  The multiplicative
coboundary $(1+a/x)^{-1}(1+b/y)$ represents the zero multiplicative
cohomology class.  Under $1+z\mapsto z$, its additive class is
\[
 \left[\frac{\epsilon+ab}{xy}\right]
   =\left[\frac{ab}{xy}\right]\ne0
       \quad\hbox{in }H^1(U,\mathcal O_U),
\]
since $\epsilon/(xy)=a/x+b/y$ is an additive coboundary whereas
$ab/(xy)$ is not.  That nonzero class reduces to zero modulo $I$.
Thus even the zero class has representatives giving different images
under the proposed canonical construction.
\end{proof}

\begin{proposition}[Proved: the comparison fails in every positive characteristic]
\label{A4:all-primes}
The restriction to characteristic two in the supplied status R2 is
incorrect.  For every prime $p$, let
\[
\begin{gathered}
 k=\mathbf F_p,\quad R=k[x,y],\quad
 A=R[a,b]/\bigl(a^p,b^p,(ya-xb)^2\bigr),\\
 I=(\epsilon),\quad\epsilon=ya-xb,\quad A'=A_x\times A_y.
\end{gathered}
\]
These noetherian rings satisfy every hypothesis of cote 128, Lemma 1,
but its canonical square-zero comparison fails.  In particular $p=3$
gives a counterexample with $2$ invertible.  More precisely, there are
independent classes $e_0,\ldots,e_{p-1}$ in $H^1(IC^\bullet)$ such that
the two full boundary images are
\[
 \delta(C_0)=\langle e_0,\ldots,e_{p-2}\rangle_k,
 \qquad
 \delta_{\log}(C_0^\times)
 =\left\{\sum_{j=0}^{p-1}c_je_j:\ c_{p-1}=c_0\right\}.
\]
\end{proposition}
\begin{proof}
The ideal $I$ is square-zero by its defining relation.  On $D(x)$ put
$t_x=a/x$; eliminating $b$ gives
\[
 A_x=R_x[t_x,\epsilon]/(t_x^p,\epsilon^2),\qquad
 b=yt_x-\epsilon/x.
\]
Indeed the additional relation $b^p=0$ is redundant, since in
characteristic $p$ it is $y^pt_x^p-\epsilon^p/x^p=0$.
Similarly on $D(y)$ put $t_y=b/y$.  On the overlap
\[
                t_x-t_y=\frac{\epsilon}{xy}.
\]
Modulo $I$, these coordinates glue to $t$ and give
$U_0\simeq U_R\times\operatorname{Spec}k[t]/(t^p)$, with the same
punctured plane $U_R$ as before.  Consequently
$C_0=R[t]/(t^p)$.  Every element is the sum of an element of $R$ and
a unit of constant term one, so $C_0=C_0^\times+\operatorname{im}B$;
elements of $R$ already belong to $B=\Gamma(U,\mathcal O_U)$.

Multiplying the transition formula by $\epsilon$ shows that the
sections $t_x^j\epsilon$ and $t_y^j\epsilon$ agree for $0\leq j<p$.
Thus $I|_U$, viewed on $U_R$, is free on
$\epsilon,t\epsilon,\ldots,t^{p-1}\epsilon$.  If
$W=R_{xy}/(R_x+R_y)$, this proves
\[
 H^1(IC^\bullet)=\bigoplus_{j=0}^{p-1}W t^j\epsilon,
 \qquad e_j=\left[\frac{t^j\epsilon}{xy}\right].
\]
The Laurent-monomial basis proves independence and $x e_j=y e_j=0$.
For $c=\sum_{j=0}^{p-1}s_jt^j$, the transition formula and $I^2=0$ give
\[
 \delta(c)=\sum_{j=1}^{p-1}j s_j(0,0)e_{j-1}.
\]
This is exactly the first displayed boundary image.  It also computes
\[
 \operatorname{im}(B\to C_0)
     =R\oplus\bigoplus_{j=1}^{p-1}(x,y)R t^j;
\]
we do not assume that $B=A$ for this more general example.

A unit of $C_0$ has nonzero constant coefficient in $k$, and its
logarithmic boundary, after division by that constant, is obtained by
evaluating its coefficients at $(0,0)$ and applying
\[
 1+t k[t]/(t^p)\longrightarrow k[t]/(t^p),
                    \qquad u\longmapsto u'/u.
\]
This homomorphism is injective: $u'=0$ for a polynomial of degree less
than $p$ forces $u$ to be constant.  Its image has $p^{p-1}$ elements.
Factor $u=\prod_i(1-\alpha_i t)$ over an algebraic closure.  The
coefficient of $t^{p-1}$ in $u'/u$ is
$-\sum_i\alpha_i^p=(-\sum_i\alpha_i)^p$, the $p$-th power of its
constant coefficient.  Both coefficients lie in $\mathbf F_p$, so
they are equal.  This hyperplane also has $p^{p-1}$ elements, proving
the second boundary formula in full.

In particular $u=1+t$ has logarithmic boundary
\[
                 e_0-e_1+e_2-\cdots+(-1)^{p-1}e_{p-1},
\]
which is not an additive boundary.  Concretely,
$(1+t_y)^{-1}(1+t_x)$ is a multiplicative coboundary; its square-zero
image is nonzero modulo the additive boundaries because its
$e_{p-1}$ coefficient is one.  This is the same well-definedness
failure as in characteristic two.

The exact certificate \texttt{code/A4\_all\_primes\_boundary.py}
enumerates the entire coefficient rings for $p=2,3,5$, verifies both
full boundary images and this witness, and checks the truncated
derivation identities.  The argument above proves the assertion for
every prime, without extrapolating from those finite enumerations.
\end{proof}

\begin{remark}[Exact scope of the refutation]\label{A4:boundary-gate}
The original scan 20 says ``un isomorphisme canonique,'' and scan 21
identifies its two kernels by the square-zero map and the corresponding
boundary images.  Proposition~\ref{A4:boundary-counterexample} refutes
that construction and the required equality of boundary images.
It does not assert that the two groups in this particular example are
nonisomorphic as abstract groups: both are vector spaces over $\mathbf F_2$
of countably infinite dimension.  The statement with no specified meaning
of ``canonical'' must not be silently strengthened into an abstract
nonisomorphism claim.

Proposition~\ref{A4:all-primes} also refutes the supplied claim that the
canonical comparison is correct in odd characteristic; the full
boundary formulas give the obstruction there as well.
The algebra $A'$ in these examples is not finite over $A$, and $A$ is not semilocal.
Consequently this example does not refute the subsequent finite semilocal
theorem on scans 26--30.  Its vanishing assertions, including arbitrary-rank
flat-module descent, are proved independently in
Theorem~\ref{A4:finite-descent-theorem}.  The exact always-valid replacement for Lemma 1 is
Proposition~\ref{A4:boundary}; equality of its two boundary subgroups is
the precise additional condition.  In particular, surjectivity of
$H^0(C^\bullet)\to C_0$ makes both boundaries zero and suffices.

The attached exact characteristic-two script verifies the products and
the Laurent coefficient obstruction used above.  Earlier finite searches
through truncated monomial and nonmonomial algebras and through closed
covers found no counterexample in their finite families; those failed
approaches and their outputs are retained in the research log.
\end{remark}

\subsection{Vanishing survives: the finite-algebra derivation argument}

\begin{lemma}[Proved: finite-dimensional logarithmic orbits]
\label{A4:derivation-orbits}
Let $k$ be a field, $D$ a finite-dimensional commutative $k$-algebra,
$H$ a finite-dimensional $D$-module, and $\delta:D\to H$ a $k$-derivation.
For $r\geq1$ let $\mathrm{GL}_r(D)$ act on $\mathbb M_r(H)$ by
\[
       g*h=g h g^{-1}+\delta(g)g^{-1}.
\]
The following conditions are equivalent: $\delta$ is surjective; this
action is transitive for one integer $r\geq1$; it is transitive for every
integer $r\geq1$.  In particular, additive vanishing and logarithmic
vanishing are equivalent here even when their boundary subgroups differ.
\end{lemma}
\begin{proof}
Write $S=\ker\delta$.  The stabilizer of zero is $\mathrm{GL}_r(S)$.
This statement also holds as a statement about affine group schemes over
$k$: kernels commute with scalar extension, and an element of $S$ which
is invertible in $D$ is invertible in $S$, since $D$ is finite and integral
over $S$.  Both groups are smooth open subschemes of the affine spaces
of their matrices.  Consequently the dimension of the zero orbit is
\[
       r^2(\dim_k D-\dim_k S)=r^2\dim_k\operatorname{im}\delta.
\]

Suppose $\delta$ surjective and first extend scalars to an algebraic
closure of $k$.  Decompose $D$ into its local factors.  Derivations kill
idempotents, so $S$ contains their scalar span, a subalgebra $k^s$.
Put $\mathfrak n=\operatorname{rad}D$.  Then
\[
 \mathrm{GL}_r(D)=
 (1+\mathbb M_r(\mathfrak n))\rtimes\prod_{i=1}^s\mathrm{GL}_r(k).
\]
The second factor fixes zero.  Hence the zero orbit is an orbit of the
unipotent group $1+\mathbb M_r(\mathfrak n)$.  It is closed by the
Kostant--Rosenlicht theorem, whose exact statement and proof were opened
in \cite[Theorem 17.64, p.~374]{MilneAG}.  It has the dimension of the
whole affine space $\mathbb M_r(H)$ and therefore equals that space.

For the descent back to the original $k$, the fibre of the orbit map
over any $h\in\mathbb M_r(H)$ is a torsor under
$\operatorname{Res}_{S/k}\mathrm{GL}_r$.  Such a torsor is a finite
locally free $S$-module of constant rank $r$, by faithfully flat module
descent.  Every such module is free because $S$ is Artinian and semilocal.
Thus the torsor has a $k$-point, proving transitivity over $k$ itself.

Conversely suppose the action is transitive on $k$-points.  If $k$ is
infinite, these points are Zariski dense in $\mathbb M_r(H)$, so the
displayed dimension equality forces $\delta$ to be surjective.
If $k$ is finite, of characteristic $p$, then $D^p\subset S$.  Frobenius
is bijective on $D_{\mathrm{red}}$, a product of finite fields, so
$S_{\mathrm{red}}\simeq D_{\mathrm{red}}$.  Reduction of invertible
matrices through a nilpotent ideal is surjective, with kernel of
cardinality the $r^2$-th power of the ideal's cardinality.  Writing
$q=|k|$, transitivity gives
\[
 q^{r^2\dim H}
 =\frac{|\mathrm{GL}_r(D)|}{|\mathrm{GL}_r(S)|}
 =q^{r^2(\dim D-\dim S)}.
\]
Again $\delta$ is surjective.  This argument includes $r=1$ and does not
replace a finite residue field by an infinite one without justification.
\end{proof}

\begin{proposition}[Proved: the finite-rank equivalences of cote 128]
\label{A4:finite-rank-equivalence}
Let $A$ be semilocal and noetherian, let $\mathfrak m=\operatorname{rad}A$,
and let $A'$ be a finite $A$-algebra.  Put
$A_n=A/\mathfrak m^{n+1}$ and $A'_n=A'\otimes_A A_n$.
For any fixed integer $r\geq1$, the following are equivalent:
\[
 \begin{split}
 &H^1(A'_n/A_n,\mathbf G_a)=0\quad\hbox{for every }n\geq0;\\
 &H^1(A'_n/A_n,\mathrm{GL}_r)=\{1\}\quad\hbox{for every }n\geq0.
 \end{split}
\]
If they hold for one $r$, they hold for all finite $r$, including $r=1$.
Thus parts (i), $(\mathrm{ii}_1)$ and $(\mathrm{ii}_m)$ of the manuscript's
Theorem (a) survive in characteristic two.
\end{proposition}
\begin{proof}
The Artinian rings $A_n$ decompose according to the finitely many maximal
ideals of $A$, so work on one local factor, with residue field $k$.
For $n=0$, a nonzero finite $k$-algebra is faithfully flat over $k$;
faithfully flat descent gives both assertions.  A zero factor contributes
only trivial cohomology and may be omitted.

For the step from $n-1$ to $n$, put
$I=\mathfrak m^n/\mathfrak m^{n+1}\subset A_n$ and use the notation of
Proposition~\ref{A4:boundary} for $A_n\to A'_n$.
Here $I^2=\mathfrak m I=0$.  Thus $H=H^1(IC^\bullet)$ is a finite
dimensional vector space over $k$.  The connecting map
$\delta:C_0\to H$ is a derivation: expanding a product of two lifts gives
the Leibniz identity, and the product of two error terms is zero because
$I^2=0$.  Its action on $H$ and its derivation both factor through
\[
                 D=C_0/\mathfrak m C_0,
\]
a finite-dimensional $k$-algebra.  Indeed $\delta$ vanishes on scalars
from $A_{n-1}$, and $\mathfrak m H=0$.

The additive kernel is $H/\delta(D)$.  For matrices the nonabelian
cohomology exact sequence identifies the kernel with the orbit space
of $\mathbb M_r(H)$ under
\[
               h\longmapsto g h g^{-1}+\delta(g)g^{-1},
               \qquad g\in\mathrm{GL}_r(C_0).
\]
This follows directly by lifting a trivialization of the reduced
cocycle, writing the resulting cocycle as $1+z$ with $z\in IC^1$, and
changing that trivialization by an invariant matrix modulo $I$.
The action factors through $\mathrm{GL}_r(D)$, and every matrix in that
group lifts, since $\mathfrak m C_0$ is nilpotent.
Lemma~\ref{A4:derivation-orbits} says exactly that these two kernels are
trivial simultaneously.  If the lower cohomology groups vanish, they
are the entire cohomology groups in question.  Induction gives the stated
equivalence, in either direction, for every finite initial range of $n$,
and hence for all $n$.
\end{proof}

\begin{lemma}[Proved: the arbitrary-rank extension]
\label{A4:infinite-orbits}
In Lemma~\ref{A4:derivation-orbits}, if $\delta$ is surjective, its
affine matrix action is transitive also for matrices on $D^{(\Lambda)}$
for an arbitrary index set $\Lambda$.  Matrices here are column-finite;
an $H$-valued matrix means a map
$k^{(\Lambda)}\to H\otimes_k k^{(\Lambda)}$.
\end{lemma}
\begin{proof}
We first work over an algebraically closed field and in the case
$S=\ker\delta=k$.  Derivations kill idempotents, so $D$ is local;
write $N=\operatorname{rad}D$.  The restriction $\delta:N\to H$ is an
isomorphism of vector spaces.  Choose a basis $h_1,\ldots,h_d$ of $H$
and put $n_i=\delta^{-1}(h_i)$.  Define operators on the finite-dimensional
space $N$ by
\[
                 R_i(n)=\delta^{-1}(n h_i).
\]
For any vector space $V$, and endomorphisms $a_i$ of $V$, put
$h=\sum_i h_i\otimes a_i$ and $T_h=\sum_i R_i\otimes a_i$ on $N\otimes V$.
The equation for a horizontal vector $v\in D\otimes V$ is
$\delta(v)=hv$.  On the subspace $N\otimes V$ it is exactly
$(1-T_h)v=0$.

If $V$ is finite-dimensional, Lemma~\ref{A4:derivation-orbits} supplies
$g$ with $\delta(g)=hg$.  All horizontal vectors are then $g(k\otimes V)$.
Reduction of $g$ modulo $N$ is invertible, so no nonzero horizontal vector
lies in $N\otimes V$.  Therefore $1-T_h$ is invertible for every finite
matrix substitution $a_1,\ldots,a_d$.

We show that this forces the operators $R_i$ to be jointly nilpotent.
If a positive word $W$ in them were not nilpotent, use a vector space
with basis indexed by the positions in that word, and let each $a_i$
be the sum of the corresponding labelled edges in a directed cycle.
Then a power of $T_h$ has $W$ as a diagonal block.  It has a nonzero
eigenvalue; after scaling all $a_i$ by one nonzero scalar, $T_h$ has
eigenvalue one.  This contradicts the preceding invertibility.
Thus every positive word is nilpotent.

For clarity, the finite-dimensional nil-semigroup argument needed here
is included.  Let $E\subset\operatorname{End}_k(N)$ be the unital algebra
generated by the $R_i$, and $J$ the ideal spanned by their positive words.
If $J$ were not contained in the Jacobson radical of $E$, some simple
$E$-module $W$ would have $JW\ne0$.  The image of $E$ on $W$ is all of
$\operatorname{End}_k(W)$.  Indeed Schur's lemma gives
$\operatorname{End}_E(W)=k$; induction on a linearly independent list
$w_1,\ldots,w_l$ proves surjectivity of
$E\to W^l$, $e\mapsto(ew_1,\ldots,ew_l)$.
In the induction, a zero kernel in the last coordinate would make the
image the graph of an $E$-linear map $W^{l-1}\to W$, a scalar linear
combination of the projections, contradicting independence.
Applying this to a basis proves the assertion about the image algebra.
The image of $J$ is consequently the whole matrix algebra.  But every
positive word acts nilpotently on $W$, because its power is already zero
in $E$.  Every element of the image of $J$ would therefore have trace zero,
contradicting the trace-one matrix unit.  This works also in positive
characteristic.  Hence $J$ is contained in the Jacobson radical, which is
nilpotent for a finite-dimensional algebra.  Thus $J^q=0$ for some $q$.
This is the required special case of the nil-semigroup theorem, with its
proof supplied rather than an unverified appeal to a general nil-algebra
claim.

It follows that $T_h^q=0$ for arbitrary $V$ and arbitrary endomorphisms
$a_i$, since each coefficient of that power is a word of length $q$ in
the $R_i$.  For each $v\in V$ define
\[
 g(v)=1\otimes v+
       (1+T_h+\cdots+T_h^{q-1})\sum_i n_i\otimes a_i(v),
\]
and extend $D$-linearly.  This is column-finite, satisfies
$\delta(g)=hg$, and is invertible: $g-1$ has entries in the nilpotent
ideal $N$, so its inverse is a finite geometric sum.  The case $S=k$
is proved for every rank.

Now allow an arbitrary finite-dimensional $S=\ker\delta$, still over
the algebraically closed field.  Decompose by its idempotents and work
on a local factor.  Put $\mathfrak r=\operatorname{rad}S$ and
$Q=\mathfrak rD$.  Since $\delta$ is $S$-linear and surjective,
$\delta(Q)=\mathfrak rH$.  The induced derivation
\[
       D/Q\longrightarrow H/\mathfrak rH
\]
is surjective with kernel $k$: its kernel is $(S+Q)/Q$, and
$S\cap Q=\mathfrak r$, since $Q$ is nilpotent.
Apply the preceding case and lift the resulting change of basis through
the nilpotent ideal $Q$.  A column-finite invertible matrix lifts through
any nilpotent ideal: lift it and its inverse, and correct their products
by finite geometric sums.  After this change of basis the remaining
matrix lies in $\mathfrak rH$.

Inductively, if it lies in $\mathfrak r^jH$, use
$\delta(\mathfrak r^jD)=\mathfrak r^jH$ to choose a column-finite matrix
$X$ in $\mathfrak r^jD$ whose derivative cancels it modulo
$\mathfrak r^{j+1}H$.  The change of basis $1+X$ is invertible.
Its nonlinear error terms lie in $\mathfrak r^{2j}H$, hence in
$\mathfrak r^{j+1}H$.  Nilpotence of $\mathfrak r$ finishes the process.

Finally descend from the algebraic closure to the original field.
Two such changes of basis differ by an automorphism of the free
$S$-module.  Faithfully flat module descent therefore produces an
$S$-module $P$ which becomes free of rank $\Lambda$ after scalar extension.
It is flat.  A flat module over an Artinian local ring is free even at
infinite rank: lift a basis modulo its maximal ideal, use nilpotent
Nakayama for surjectivity, and use flatness followed by nilpotent
Nakayama on the kernel for injectivity.  Apply this on the finitely many
local factors of $S$; the ranks are detected after field extension and
are all $\Lambda$.  Thus $P$ is free, and a basis gives the required
change of basis over the original field.
\end{proof}

\begin{theorem}[Proved: the finite semilocal theorem of cote 128 survives]
\label{A4:finite-descent-theorem}
Keep the notation and hypotheses of
Proposition~\ref{A4:finite-rank-equivalence}, and put
\[
 B^{(n)}=H^0(A'_n/A_n),\qquad B=H^0(A'/A).
\]
The equivalent conditions in that proposition are also equivalent to:
for every $n$, $\operatorname{Spec}A'_n\to\operatorname{Spec}B^{(n)}$
is an effective descent morphism for all flat quasi-coherent modules.
Under these conditions,
\[
 H^1(A'/A,\mathbf G_a)=0,\qquad
 H^1(A'/A,\mathrm{GL}_r)=\{1\}\quad(r\geq1),
\]
and $\operatorname{Spec}A'\to\operatorname{Spec}B$ is an effective
descent morphism for finite flat modules.
If $A$ is Artinian, then for every flat $B$-scheme $X$, the base change
$X\times_B\operatorname{Spec}A'\to X$ is an effective descent morphism
for all flat quasi-coherent modules.
These are exactly parts (a), (b), and (c) of the theorem on scans 26--30;
the proof uses vanishing, not the false canonical comparison.
\end{theorem}
\begin{proof}
\emph{(a), arbitrary rank and effectiveness.}
For a free module indexed by $\Lambda$, the endomorphism cochain complex
is the column-finite matrix complex.  Direct sums and products are exact
in the category of modules, so its degree-one cohomology is the
column-finite matrix module on $H^1(IC^\bullet)$.  Thus the same
square-zero argument as in Proposition~\ref{A4:finite-rank-equivalence},
now using Lemma~\ref{A4:infinite-orbits}, proves triviality of descent
cocycles in every free rank over every $A'_n$.

We spell out the passage to flat modules and to the equalizer base.
If $b\in B^{(n)}$, its two images in $A'_n\otimes_{A_n}A'_n$ agree.
Hence the natural map
\[
 A'_n\otimes_{A_n}A'_n\longrightarrow
 A'_n\otimes_{B^{(n)}}A'_n
\]
is an isomorphism, as are the higher analogous tensor products.
The descent data are therefore unchanged by replacing $A_n$ by
$B^{(n)}$.  This latter ring is Artinian, being a finite submodule of
$A'_n$ over $A_n$, and it embeds in $A'_n$.
On a local factor of $A_n$ with residue field $k$ and $A'_n\ne0$,
the ring $B^{(n)}$ is local with residue field $k$.  Indeed its image
modulo $\mathfrak mA'_n$ is contained in the equalizer $k$ for the
faithfully flat $k$-algebra $A'_n/\mathfrak mA'_n$, and contains that
$k$.  The kernel is nilpotent.  Thus its local-factor decomposition is
the supported part of the local-factor decomposition already used in
the induction.
On each local factor of $B^{(n)}$, the ranks of a flat $A'_n$-module with
descent datum are the same at all maximal ideals of $A'_n$: tensoring
two of their residue fields over the residue field of $B^{(n)}$ gives
a nonzero ring, and the descent isomorphism identifies the two ranks
there.  Flat modules over the Artinian local factors of $A'_n$ are free
by the argument at the end of Lemma~\ref{A4:infinite-orbits}.
Thus they are free of one rank on each factor of $B^{(n)}$, and the
arbitrary-rank triviality just proved makes their descent effective.
Full faithfulness follows coefficient by coefficient from the equalizer
definition of $B^{(n)}$, using column-finite matrices for free modules.

Conversely, an effective datum on $(A'_n)^r$ descends to a flat
$B^{(n)}$-module $P$.  It is finite: generators after tensoring with
$A'_n$ involve finitely many elements of $P$, and the remaining quotient
is zero by Lemma~\ref{A4:artinianflat}.  Its ranks on the local factors
are $r$, so it is free of rank $r$.  This proves $(\mathrm{iii})\Rightarrow
(\mathrm{ii}_r)$ and completes all equivalences in (a).

\emph{(b), the additive assertion.}
Let $z\in C^1$ be an additive cocycle, where $C^q=(A')^{\otimes_A(q+1)}$.
Its reduction modulo $\mathfrak m^{n+1}$ is a coboundary for every $n$,
so
\[
 z\in\operatorname{im}(d:C^0\to C^1)+\mathfrak m^{n+1}C^1
       \quad\hbox{for all }n.
\]
Krull intersection for the finite $A$-module
$C^1/\operatorname{im}d$, with $\mathfrak m\subset\operatorname{rad}A$,
gives $z\in\operatorname{im}d$.  Thus $H^1(A'/A,\mathbf G_a)=0$.

\emph{(b), matrices over the completion.}
First suppose $A$ complete for its radical topology.  Fix a matrix
cocycle, and let $T_n$ be its nonempty set of trivializations over $A'_n$.
It is a torsor under $\mathrm{GL}_r(B^{(n)})$.
For fixed $n$, the images of $B^{(N)}\to B^{(n)}$, $N\geq n$, form a
descending family of submodules of a finite Artinian module, and hence
stabilize.  The kernel of each such ring map is nilpotent: it is
contained in $\mathfrak m^{n+1}A'_N$.
Consequently the image of the corresponding general linear group is
the general linear group of the image ring.  Here invertibility in that
image ring is detected in $B^{(n)}$, which is finite and integral over it.
Thus the images of these groups stabilize as well.  The images
$T_N\to T_n$ are nonempty nested cosets of those groups and consequently
stabilize.  The countable inverse system $(T_n)$ is Mittag--Leffler with
nonempty terms, so its limit is nonempty.  Since
$A'=\varprojlim A'_n$, a compatible system gives an invertible matrix
trivializing the original cocycle.  This proves the matrix assertion
over a complete base, without assuming that individual transition maps
$B^{(n+1)}\to B^{(n)}$ are surjective.

\emph{(b), descent from completion and finite flat modules.}
Completion $A\to\widehat A$ is faithfully flat.  The ring
$\widehat B=B\otimes_A\widehat A$ is the equalizer after completion,
because all the modules in that equalizer are finite and completion is
flat.  For a finite flat $A'$-module $E'$ with datum, its invariant
module $E$ is finite over $A$ and satisfies
$E\otimes_A\widehat A=H^0_D(\widehat A'/\widehat A,\widehat E')$.
Over each local factor of $\widehat A$ the equalizer $\widehat B$ is
local with the same residue field: modulo the radical of $\widehat A$,
invariant elements lie in the base field by faithfully flat descent,
and integrality detects units.  The descent datum gives equal ranks at
the maximal ideals of $\widehat A'$ as in (a).  Hence $\widehat E'$ is
free on that factor, and the complete-base matrix assertion makes its
datum descend to a finite free $\widehat B$-module.
Faithfully flat descent of flatness and of isomorphisms now shows that
$E$ is finite flat over $B$ and that $E\otimes_B A'\to E'$ is an
isomorphism.  Full faithfulness again follows from the equalizer,
using local finite free presentations.  Applied to $(A')^r$, the
descended module is finite projective of constant rank $r$ over the
semilocal ring $B$, hence free.  This proves the remaining assertions
of (b).

\emph{(c), arbitrary flat base change.}
When $A$ is Artinian, (a) gives effectiveness for all flat modules over
$B\hookrightarrow A'$.  Work on an affine open
$\operatorname{Spec}T$ of the flat $B$-scheme $X$ and put
$T'=T\otimes_B A'$.  A flat $T'$-module with descent datum is flat as an
$A'$-module, since $T'$ is flat over $A'$.  Forgetting its $T$-action,
use (a) to descend it to a flat $B$-module $E$.  Full faithfulness descends
the multiplication operators from $T$, with their additive and
multiplicative relations, so $E$ is a $T$-module and
$E\otimes_B A'$ is the given $T'$-module.

It remains to check $T$-flatness, which must not be inferred merely from
injectivity of $T\to T'$.  Work on an Artinian local factor $(B,\mathfrak b,k)$
and choose a residue field $l$ of $A'$ above it.  The given flatness,
followed by base change to $l$, and faithfully flat descent along $k\to l$,
implies that $E\otimes_B k$ is flat over $T\otimes_B k$.
Both $T$ and $E$ are flat over $B$.  A free $B$-resolution of $k$,
tensored with $T$ and then with $E$, gives
\[
 \operatorname{Tor}_i^T(T/\mathfrak bT,E)
       =\operatorname{Tor}_i^B(k,E)=0\quad(i>0).
\]
For every $T/\mathfrak bT$-module $M$, change of rings therefore gives
$\operatorname{Tor}_i^T(M,E)=
\operatorname{Tor}_i^{T/\mathfrak bT}(M,E/\mathfrak bE)=0$ for $i>0$.
Filter an arbitrary $T$-module by its finitely many powers of the
nilpotent ideal $\mathfrak b$; the Tor exact sequences give
$\operatorname{Tor}_1^T(M,E)=0$ for every $M$.  Thus $E$ is $T$-flat.
The same underlying-module argument gives full faithfulness, including
$T$-linearity.  These constructions agree on affine overlaps by
uniqueness and glue to $X$, proving (c).
\end{proof}

\subsection{Projectivity over a normal locally noetherian base}

\begin{theorem}[Known, with its projectivity consequence proved]
\label{A5:projectivity}
Let $f:A\to S$ be an abelian scheme and let $S$ be normal and locally
noetherian.  There exists an $f$-ample invertible sheaf $L$ on $A$.
In fact $A$ admits a polarization $\lambda:A\to A^\vee$, and one may take
$L=(1,\lambda)^*\mathcal P$, where $\mathcal P$ is the normalized Poincar\'e
bundle on $A\times_S A^\vee$.
The sheaf $V=f_*(L^{\otimes3})$ is finite locally free, its formation
commutes with base change, and evaluation gives a closed immersion
\[
             A\lhook\joinrel\longrightarrow\mathbf P_S(V).
\]
Here $\mathbf P_S(V)$ parametrizes invertible quotients of $V$.
Thus the assertion on cote 131, scan 17, is correct in the relative
ample-bundle, or projective-bundle, sense of projectivity.
\end{theorem}
\begin{proof}
Conrad's Lemma 3.2.1 \cite[p.~18]{ConradCanonical} is precisely the
polarization assertion, with the stated normal locally noetherian
hypotheses; we opened both statement and proof.  Its construction works
on the open connected components of $S$.  On each such integral component
a generic polarization extends uniquely as a homomorphism
$A\to A^\vee$ by the Weil extension theorem.  The proof then checks that
$(1,\lambda)^*\mathcal P$ is relatively ample, by openness of the ample
locus and specialization over discrete valuation rings.  The line bundles
on the disjoint components give a line bundle on all of $A$.  In particular,
quasi-compactness of $S$ is not an additional hypothesis in the cited lemma.

For completeness, the passage from that result to an embedding is as
follows.  On every geometric fibre, the cube of an ample line bundle on
an abelian variety is very ample \cite[Remark 7.3]{MilneAV}.
Its higher cohomology vanishes: the index theorem
\cite[Theorem 13.3(c)]{MilneAV} gives cohomology in exactly one degree,
and very ampleness gives a nonzero space of sections in degree zero.
Cohomology and base change for the proper flat morphism $f$ now gives
the asserted locally free sheaf $V$ and base-change comparison.
The evaluation map $f^*V\to L^{\otimes3}$ is surjective, as follows
on fibres and then by Nakayama's lemma.  It defines the displayed map;
it is a closed immersion because that property holds on every fibre for
this proper morphism.  All these assertions are local on $S$, so they
remain valid when the ranks on its components are unbounded.
\end{proof}

\begin{remark}[Proved: the convention about projectivity matters]
\label{A5:fixed-projective-space}
If ``projective'' is defined to require a closed immersion in one fixed
$\mathbf P^N_S$, the statement for arbitrary locally noetherian $S$ needs
that convention distinguished.  For example, take an elliptic curve $E/k$,
$S=\coprod_{r\geq1}\operatorname{Spec}k$, and
$A=\coprod_{r\geq1}E^r$.  This is an abelian scheme over a normal locally
noetherian base and has the relatively ample bundle of
Theorem~\ref{A5:projectivity}; no fixed $\mathbf P^N_S$ can contain it,
since the fibre dimensions $r$ are unbounded.  The theorem supplies the
precise global conclusion $A\hookrightarrow\mathbf P_S(V)$, as well as
embeddings in ordinary projective space locally on the base.
\end{remark}

\begin{proposition}[Proved: what the correspondence argument requires]
\label{A5:correspondences}
The scheme of divisorial correspondences packages the polarization in
Theorem~\ref{A5:projectivity}, but its separated locally quasi-finite
representability alone does not prove projectivity.  The needed additional
inputs are extension of the generic homomorphism over the normal base
and preservation of its polarization property under specialization.
With these inputs the argument on cote 131, scan 17, works.
\end{proposition}
\begin{proof}
A divisorial correspondence on $A\times_S A$ rigidified along both zero
sections determines a homomorphism $A\to A^\vee$; conversely the Poincar\'e
bundle pulls back to that correspondence.  Thus the generic polarization
is a generic section of this correspondence functor.  The extension
theorem used in Conrad's proof supplies its unique global section, and
the positivity step in the same proof makes its diagonal pullback ample.
These are precisely the two operations required to pass from the
correspondence to a relatively ample bundle.

For the insufficiency assertion, let $S=\operatorname{Spec}k[t]$ and
$U=D(t)$.  The map $U\to S$ is separated and locally quasi-finite, and
has a section over the generic point of $S$.  That section cannot extend
to $S$, since such a section would put the point $t=0$ in the image of
$U$.  Consequently the representability property cited on the manuscript
page cannot, by itself, perform the generic-section extension.  Normality
enters through the extension theorem for the actual abelian schemes, not
through a general extension property of locally quasi-finite schemes.
\end{proof}

